%%  Build with:  latexmk -lualatex writeup.tex
\documentclass{article}

\usepackage{ebproof}
\usepackage{listings}
\usepackage{xcolor}
\usepackage{geometry}

\geometry{margin=1.5in}

\lstset{
  frame=none,
  xleftmargin=2pt,
  stepnumber=1,
  numbers=left,
  numbersep=5pt,
  numberstyle=\ttfamily\tiny\color[gray]{0.3},
  belowcaptionskip=\bigskipamount,
  captionpos=b,
  escapeinside={*'}{'*},
  language=haskell,
  tabsize=2,
  emphstyle={\bf},
  commentstyle=\it,
  stringstyle=\mdseries\rmfamily,
  showspaces=false,
  keywordstyle=\bfseries\rmfamily,
  columns=flexible,
  basicstyle=\small\sffamily,
  showstringspaces=false,
  morecomment=[l]\%,
}

\newcommand{\todo}[1]{\textcolor{red}{#1}}

\title{Featherweight ATS3 Report}
\author{Thomas Ulmer}

\begin{document}
\maketitle

\section{Overview}

This document presents FWATS3, a stripped back version of the
programming language ATS3 implemented in order to formalize the
semantics of templates. The core language over which templates are
applied is intentionally underspecified and primitive. We give a
description of the languge, some machinery for type level
substitutions, templates and the rules that govern them, alternatives,
and extensions to the current system.

\section{Core language}

The core of FWATS is a small calculus with only a few primitive
concepts. It features simple types, with n-ary tuples and function
types as the only non-builtin constructors. Multi-argument functions
are constructed with a single tuple argument and the use of the lambda
or fix rule, depending on the need for recursion. Application, let
binding, and variable references follow the standard lambda calculus
constructions, and the sole conditional \verb|if| is similarly
simple. We will avoid giving concrete semantics for this language, as
discussed later. \todo{where?}

The syntax is split between expressions and declarations and
is summarized as follows:

\begin{lstlisting}[language=Haskell]
data SExp n
  = SCon n [(SExp n)]
  | SFun (SExp n) (SExp n)
  | STuple [(SExp n)]

data DExp n
  = DInt Int
  | DBool Bool
  | DStr String
  | DOp1 n (DExp n)
  | DOp2 n (DExp n) (DExp n)
  | DVar n
  | DLam (n, (SExp n)) (DExp n)
  | DFix n (n, (SExp n)) (DExp n) (SExp n) -- self name, arg+type, body, overall type
  | DApp (DExp n) (DExp n)
  | DIf (DExp n) (DExp n) (DExp n)
  | DTuple [(DExp n)]
  | DProj Int (DExp n)
  | DLet [DDecl n] (Maybe (DExp n))

data DDecl n
  = DBind n (DExp n)
  | DLocal [(DDecl n)] [(DDecl n)]

data DProgram n
  = DProgram [(DDecl n)] (DExp n)
\end{lstlisting}

Here \verb|n| is the type that names range over. The convention is
that \verb|DProgram| and \verb|DLet| declarations range over both
subsequent declarations and the body. The construction \verb|DLocal|
allows for scoped definitions that range over other declarations. The
effect is similar to having the right side of each exported
declaration be a let binding, but allows sharing of the private or
local bindings across several definitions. Here again the convention
is sequential binding ranging downward to the end of the scope,
allowing for iterative definitions. Thus the only source of cyclic
references is \verb|DFix|.

\subsection{Core types}

This language contains only primitive type constructors. The only
non-nullary constructors are \verb|STuple| and \verb|SFun| which are
simple product types and simple function types. Type checking this
language is trivial and can proceed entirely with bottom up type
construction. Explicit annotation could be required on binds, but is
elided here. 

We will take the notational convention of representing expression with
their Haskell AST counterparts offset in square brackets. \todo{Is
  there a better way?}

\todo{one rule for each primitive type}
\todo{one rule for operators}

\begin{displaymath}
  \begin{prooftree}
    \hypo{\Gamma , a : \alpha \vdash b : \beta}
    \infer1[lam]{ \Gamma \vdash [ $DLam$\: (a,\: \alpha)\: b ] : 
      [$SFun$ \alpha \: \beta] }
  \end{prooftree}
\end{displaymath}

\begin{displaymath}
  \begin{prooftree}
    \hypo{\Gamma , s : [$SFun$ \: \alpha \: \beta ] , a : \alpha \vdash b : \beta}
    \infer1[fix]{ \Gamma \vdash [$DFix$ \: s \: (a,\: \alpha) \: b \:
      [$SFun$ \: \alpha \: \beta]]) : [$SFun$ \: \alpha \:\beta ] }
  \end{prooftree}
\end{displaymath}

\begin{displaymath}
  \begin{prooftree}
    \hypo{\Gamma \vdash c : Bool }
    \hypo{\Gamma \vdash x : \alpha }
    \hypo{\Gamma \vdash y : \alpha }
    \infer3[if]{ \Gamma \vdash [$DIf$ \: c \: x \: y] : \alpha }
  \end{prooftree}
\end{displaymath}

\begin{displaymath}
  \begin{prooftree}
    \hypo{\Gamma \vdash f : [$SFun$ \: \alpha \: \beta ] }
    \hypo{\Gamma \vdash a : \alpha  }
    \infer2[app]{ \Gamma \vdash [$DApp$ \: f \: a] : \beta }
  \end{prooftree}
\end{displaymath}


\todo{This isn't careful enough, it's abuse of notation that v is a
  name here but would be a DExp in other rules. Be clear what the
  context looks like}
\begin{displaymath}
  \begin{prooftree}
    \hypo{\Gamma \vdash v : \alpha }
    \infer1[var]{ \Gamma \vdash [$DVar$ \: v] : \alpha }
  \end{prooftree}
\end{displaymath}

\todo{one rule for tuples}
\todo{one rule for projection}
\todo{one rule for let}

\section{Substitutions}

In order to talk about templates ranging over more than a single type,
we need machinery to handle matching partially constrained types and
deriving and applying substitutions. We take said machinery from
``Typing Haskell in Haskell'' \todo{cite}, which presents a
comprehensible view of structural matching producing substitutions,
different notions of composing or combining substitutions, and
applying substitutions to types containing type variables. Their
notions of classes and constraints are not relevant here. 

Specifically we define \verb|Subst n| to be a finite map of names to
\verb|SExp n|, i.e. types which may contain further type variables. We
then define left sided matching like so:
\begin{lstlisting}
sMatch :: (SExp Name) -> (SExp Name) -> TM s r (Subst Name)
sMatch (SVar lv) r =
  pure $ substAtom lv r
sMatch l@(SCon ln las) r@(SCon rn ras)
  | (ln == rn) =
    let folder ctx (l, r) = sMerge ctx =<< sMatch l r
    in foldM folder substEmpty (zip las ras)
  | otherwise = gErr $ "Can't match " ++ show l ++ " with " ++ show r
sMatch (SFun lf la) (SFun rf ra) = do
  m1 <- sMatch lf rf
  m2 <- sMatch la ra
  sMerge m1 m2
sMatch (STuple ls) (STuple rs) =
  let folder ctx (l, r) = sMerge ctx =<< sMatch l r
  in foldM folder substEmpty (zip ls rs)
sMatch l r = gErr $ "Can't match " ++ show l ++ " with " ++ show r
\end{lstlisting}
The details of the \verb|TM s r| monad are not relevant here, we
simply use it to thread errors. We also need the merging of two
substitutions, checking that they agree on overlap:
\begin{lstlisting}
sMerge :: (Subst Name) -> (Subst Name) -> TM s r (Subst Name)
sMerge ls rs = do
  let sint = M.intersectionWith (\a b -> (a,b)) ls rs
  let agree = foldr (\(a,b) acc -> acc && a == b) True sint
  () <- if agree
       then pure ()
       else gErr $ "Disagreement in static merge " ++ show ls ++ 
         " and " ++ show rs ++ " overlap " ++ show sint
  pure $ M.union ls rs
\end{lstlisting}
We also can apply substitutions to types:
\begin{lstlisting}
sSubst :: (SExp Name) -> (Subst Name) -> TM s r (SExp Name)
sSubst v@(SVar n) sctx =
  case M.lookup n sctx of
    Just s -> pure s
    Nothing -> pure v
sSubst (SCon n ss) ctx = do
  ss' <- mapM (\s -> sSubst s ctx) ss
  pure (SCon n ss')
sSubst (SFun f a) ctx = do
  f' <- sSubst f ctx
  a' <- sSubst a ctx
  pure $ SFun f' a'
sSubst (STuple ss) ctx =
  STuple <$> mapM (\s -> sSubst s ctx) ss
\end{lstlisting}
as well as for runtime expressions and declarations in a similar
recursive style (elided for space):
\begin{lstlisting}
sSubstD :: (DExp Name) -> (Subst Name) -> TM s r (DExp Name)
sSubstDecl :: DDecl Name -> Subst Name -> TM s r (DDecl Name)
\end{lstlisting}

\section{Templates}

We extend our core language with four additions to the core syntax
tree:
\begin{itemize}
\item A new constructor \verb|SVar n| for \verb|SExp n|.
\item Two new constructors \verb|DTDec (n, [n], SExp n)| and \\
  \verb|DTImpl (n, [n], DExp n, [SExp n], Int)| for \verb|DDecl n|.
\item A new constructor \verb|DTApp n [SExp n]| for \verb|DExp n|.
\end{itemize}

Unsurprisingly, these correspond to type variables, template
declarations and implementations, and template usage.

The template declarations serve to give a unified shape to all
implementations of the same named template. In short, it associates a
single polymorphic type to the template name that all implementations
must match with.

A template implementation is composed of a previously declared name,
some number of fresh type variables, an expression in which said type
variables may appear, a sequence of types ranging over the fresh
variables which categorize the space of applications implemented by
the body, and an integer tag that is uniquely identifies this
implementation.

A implementation conforms to its relative declaration if
\begin{itemize}
\item The names are equal.
\item The number of argument types matches the number of parameters of
  the template.
\item There is some most general substitution $s$ formed by mapping
  the declared parameter variables to the corresponding arguments,
  where all pairings agree at every overlapping point.
\item $s$ applied to the declared polymorphic type is equal to the
  type of the body of the implementation. 
\end{itemize}
In short, the implementation must represent some subset of the
polymorphic type associated with the declaration. In particular, if
the implementation has no fresh type variables that it ranges over,
then it defines some single type in the space of the polymorphic type.

Templates are the primary motivation for including \verb|DLocal|,
since declarations and implementations require more careful scoping
than dynamic binds. The primary issue is that once an implementation
is visible in the context, it can shadow part or all of the desired
implementation and there is no way to specify the prior one without
adjusting the declaration and every other associated instance. See the
subsection about expansion. \todo{reference}

\subsection{Type Checking}

We perform type checking both before and after exhaustively expanding
templates. For checking before expansion, the check on template
declarations merely checks that that the output type is sane relative
to the argument type variables. Implementations are checked for the
conformity conditions presented in the previous subsection
\todo{reference}, namely that they represent subsets of the declared
polymorphic type and agree in type level arity. Template use sites
(instances of \verb|DTApp|) are checked using only the
declaration. They must agree in number of type level arguments and are
given an overall type formed by matching the declaration's arguments
against the supplied arguments pairwise, merging (in the specific
sense defined in the section on substitutions) into a single
substitution, and apply said substitution to the result type of the
declaration. 

Critically, type checking before expansion ensures the program is well
typed without requiring any details about the mechanism by which
implementations are selected. This is the strongest semantic
statement that can be made about the program's correctness without
fully unfolding the templates, i.e. running expansion in its
entirety. \todo{Do I still believe this? Justify} 

After expansion, type checking again is not strictly required, as
templates are essentially hygienic macros. In a language with more
complicated types it may be required or desired for more
comprehensible error messages or expose further type directed
optimizations. 

\subsection{Implementation Selection}

The procedure by which a suitable implementation is selected for a
given use site is called a selection scheme. To define a scheme one
must give rules for determining the following for each use site:
\begin{itemize}
\item The context, the set of implementations to consider
\item A notion of matching, which implementations are suitable
\item A process of instantiating, how to transform the selected
  implementation for the specific use site
\end{itemize}

In this presentation and associated implementation, we present a
simple if somewhat limiting scheme: capturing first fit. This scheme
uses simple snapshots of the accumulated context to act as closures
for recursive expansion and uses the matching and substitution
machinery for matching and instantiating implementations.

\subsubsection{Capturing First Fit}

In Capturing First Fit, when a template use is encountered, we scan
backward in the template context and select the first implementation
that matches with the type level arguments of the use site. We call
this scheme ``capturing'' because when recursively expanding the body
of the template, we do so in the template context as it existed when
the selected implementation was defined. In short a template
implementation forms a closure over the template context. This allows
the consistent use of a libraries' internal templates, even in the
presence of recursive expansion outside of the context of the library.

However this scheme is somewhat limiting and can act in unintuitive
ways. For instance two adjacent templates can be invoked by the user,
with the knowledge that they are implemented using the same template
declaration internally. However despite being invoked in the same
context, they may internally use different implementations. \todo{example}

\subsubsection{Non-capturing First Fit}

Equally intuitive, we could instead consider a close relative,
Non-capturing First Fit. In this scheme the selection criteria remains
the same, but recursive body expansion happens in the same context as
the outer template. This has the benefit that it allows templates to
express some end-user extensibility. A library can give an
implementation without specifying all of the implementations it
depends on, allowing extremely general definitions in some
cases. \todo{example, for instance fast exponentiation based on
  supplied multiply}

This gives intuitive behavior in the default case, as adjacent
invocations of templates have access to the same context, specifically
in regard to user-provided dependencies. Libraries can prevent
shadowing the library-integral definitions by including the
implementation in a \verb|DLocal| inside the expanded body to ensure
it is found first. 

\subsubsection{Non-capturing Best Fit}

If the scheme's notion of matching forms a metric rather than a simple
boolean check, then one can embed other semantic program information
in the data associated with each implementation, making template
resolution serve as a search problem for the best implementation among
several that are assumed to function for the given task. This is
increasingly easy to imagine in languages with more expressive types
than the simple skeleton presented here. \todo{This feels like a cop
  out but I don't actually have that much to say}

Even without requiring a notion of a metric on matchings, one can
consider encoding a very basic permission system through the use of a
parameterized data type a la Peano numerals. An extra type argument
can be appended and implementations can match on $S(x)$ for some fresh
variable $x$ for requiring one level, or $S(S(x))$ for two, et
cetera. This allows specifying a requirement for an implementation
surpassing some security requirement, but not the desire to always
take the most or the least elevated implementation beyond some point.

\subsection{Expansion}

Expansion is largely determined by choice of selection scheme. All
that remains is walking the AST, accumulating a context of the visible
templates. Upon encountering a template use, select an implementation
according to the chosen scheme and replace the use site with body of
the implementation after applying the substitution derived from the
types given by the use site. \todo{Be more formal/mechanical about
  this process} Then continue expansion into the newly inserted
body. The choice of context during body expansion is determined by the
scheme, as described above.

In order to handle the scoping limits supplied by \verb|DLocal| the
current implementation performs some local program transformations:
First eliminating ambiguity around names by including integer tags as
part of the name during an initial binding pass. This allows us to
widen the scope of any \verb|DBind| without fear that it will change
what any \verb|DVar| points to, as the integer tags will ensure that
colliding user-supplied names don't actually compare equal. Once this
is done, we can safely promote any \verb|DBind| out of the private or
scoped half of a \verb|DLocal| either upwards into the surrounding
scope, or downwards into the public half of the same pair.

Specifically a \verb|DBind| that is referenced inside a template
implementation may need to be reachable long after the end of the
\verb|DLocal| block, particularly in the presence of recursive
expansion. With the expansion strategy of effectively inlining the
body at use sites, the simplest way to ensure that the \verb|DBind| is
visible anywhere the implementation could be reachable is simply to
promote it to the surrounding binding context, relying on name tags
not to change semantics. This transformation must respect the
definition order of definitions, and template definitions or
implementations cannot be promoted with the same impunity.

We apply the following rules that preserve the definition order but
give us the desired scoping changes when encountering \verb|DLocal| nodes:
\begin{itemize}
\item Any prefix of \verb|DBind| binds in the private block are moved
  upwards to sit immediately prior to the \verb|DLocal| block in the
  surrounding scope.
\item Any contiguous stretch of \verb|DTDec| or \verb|DImpl|
  definitions remain in the private block.
\item Let $s$ be the maximal suffix of the private block that starts with a
  \verb|DBind| but may contain all four kinds of definitions after the
  first two rules. The original body of the \verb|DLocal|, $bs$ is replaced
  with $[$DLocal$ \:s\: bs ]$ and the rules are applied to this inner
  block recursively.
\end{itemize}
\todo{Give a concrete example}

This transformation can be applied in the same pass as general
template expansion.

\todo{Talk about other ways to solve this problem, think scoping w/ a
  visible tag set for each impl rather than explicit scopes, then drop
public vs private entirely, more flexible I think.}

\todo{What else?}

\section{Extensions}

The system described here can be extended in a number of ways. The
most obvious limitation is presented currently is that the underlying
core language is quite limited. A more reasonable language with sum
types, matching, and more built-in types could be employed without
having to make any changes to the template system whatsoever.

\subsection{Higher Kinded Type Variables}

More interestingly, if type variables are allowed to range over type
constructors, then this system recovers a lot of the flavor of Haskell
type classes, while being exhaustively expanded at compile time. For
instance on could write accumulations over arbitrary containers with
the use of a \verb|foldr| template, or arbitrary monads in the
standard Haskell style. The substitution handling is already
expressive enough to handle this, though it lacks some safety checks
around matching kinds at present.

\subsection{Runtime Polymorphism}

If the core language supports functions that are truly polymorphic, i.e.
a single function body's code that can be executed for any instance
rather than the templates single source body expanded to many runtime
functions, then powerful generic implementations of some templates may
be possible. This is the standard practice of ML and its language
descendants, where all polymorphic functions have a single generic
implementations and many compiler generated monomorphic variants for
performance. By exposing both polymorphic functions and templates to
the user, the language provides much more direct control and
eliminated guesswork and cajoling the compiler to optimize or inline
in the desired way.

Again we benefit from Haskell's prior work, as the substitution
machinery we inherit from them is already suited to work over
polymorphic types. It may be useful to represent type variables
introduced by templates which should be exhaustively eliminated and
type variables associated with polymorphic functions with separate
constructors or a flag so as to aid debugging, but again it is not
strictly necessary. \todo{Is this true? I think so but I'm not
  completely sure}

\subsection{Separate Compilation}

Templates present a serious barrier to efficient separate
compilation. If templates never stretch across module or compilation
boundaries then no changes are necessary, but this is a serious
limitation given the expressive power of templates.

Templates in this presentation are fundamentally source to source
transformations. As such, any attempt to delay resolution requires
delaying a large section of the compiler's work until the linking
phase (as compiling any C++ project would have you know). Depending on
the exact details of the selection scheme and intra-module visibility,
it may be possible to do meaningful work without a complete program
view, for instance partially evaluating template selection in a
capturing scheme, where some recursive implementation selections are
decidable without full information. In any case this requires
infrastructure for compiled headers or interface files or other
details of the compilation process that depend critically on the
details of the core language and are out of scope here.

%% \printbibliography
\end{document}
